\documentclass{article}
\usepackage{amsmath, amssymb, amsthm}
\usepackage{hyperref}
\usepackage{geometry}
\geometry{margin=1in}

\title{Erd\H{o}s Problem \#517}
\date{Last updated: 2025-08-31}
\author{Source: erdosproblems.com}

\begin{document}
\maketitle

\noindent\textbf{Status:} OPEN \quad \textbf{No prize}

\noindent\textbf{Tags:} analysis

\medskip
\noindent\textbf{Problem Statement:}

Let $f(z)=\sum_{k=1}^\infty a_kz^{n_k}$ be an entire function (with $a_k\neq 0$ for all $k\geq 1$). Is it true that if $n_k/k\to \infty$ then $f(z)$ assumes every value infinitely often?

\bigskip
\noindent\textbf{Remarks:}

A conjecture of Fej\'{e}r and P\'{o}lya. 

Fej\'{e}r \cite{Fe08} proved that if $\sum\frac{1}{n_k}<\infty$ then $f(z)$ assumes every value at least once, and Biernacki \cite{Bi28} proved that if $\sum\frac{1}{n_k}<\infty$ then $f(z)$ assumes every value infinitely often.

P\'{o}lya \cite{Po29} proved that if $f$ has finite order then $f(z)$ assumes every value infinitely often under the assumption that $\limsup (n_{k+1}-n_k)=\infty$.

\bigskip
\noindent\textbf{References:}

[Bi28] Biernacki, Mi\'{e}cislas, Sur les \'{e}quations alg\'{e}briques contenant des param\'{e}tres arbitraires. (1928), 145.
 
 [Fe08] Fej\'{e}r, Leopold, \"{U}ber die Wurzel vom kleinsten absoluten Betrage einer algebraischen Gleichung. Math. Ann. (1908), 413-423.
 
 [Po29] P\'olya, G., Untersuchungen \"uber {L}\"ucken und Singularit\"{a}ten von
{P}otenzreihen. Math. Z. (1929), 549--640.

\bigskip
\noindent\small{Source: \url{https://www.erdosproblems.com/517}}
\end{document}
